\vspace{-6mm}
\section{Introduction}
\label{sec:introduction}

Owing to its non-ionizing, real-time, portable, and low-cost nature, ultrasound
(US) is widely used in diagnosis, monitoring, and image-guided interventions
\cite{tsedendambaReviewLinearArray2025}.
However, routine US remains predominantly 2D, requiring sonographers to sweep
regions of interest and mentally integrate sequential cross-sections into 3D
anatomy, increasing visuospatial burden and operator dependence
\cite{tenajasRecentAdvancesAI2023}.
Although dedicated 3D/4D probes directly acquire volumetric data, broader
adoption is limited by specialized hardware and trade-offs in field of view,
spatial and temporal resolution, and system complexity
\cite{ingramThreeDimensionalUltrasound2025}.
Consequently, reconstructing 3D US volumes from routine 2D B-mode sequences
offers a flexible and cost-effective alternative for volumetric visualization,
quantitative assessment, and downstream clinical analysis.
To characterize this paradigm, we consider two complementary aspects:
acquisition setting and reconstruction pipeline design.

From the acquisition perspective, methods are broadly categorized as tracked or
sensorless according to how frame-wise spatial information is obtained
\cite{adriaansTrackerless3DFreehand2024}.
Tracked reconstruction synchronizes B-mode frames with probe poses measured by
optical or electromagnetic devices, requiring extra hardware, calibration, and
setup that increase cost and hinder routine deployment
\cite{mozaffariFreehand3DUltrasound2017,
adriaansTrackerless3DFreehand2024}.
By contrast, sensorless reconstruction estimates inter-frame geometry directly
from US sequences, eliminating dedicated tracking hardware and improving
portability and accessibility, and has therefore attracted increasing attention
in freehand 3D US reconstruction
\cite{adriaansTrackerless3DFreehand2024,liTUSREC2024}.

From the pipeline perspective, 3D US reconstruction involves geometry
localization and scene representation.
Geometry localization estimates inter-frame transformations before localized
slices are assembled by interpolation or compounding, making pose estimation
the primary focus
\cite{prevostDeepLearningSensorless2017,guoSensorlessFreehand3D2020,
guoUltrasoundVolumeReconstruction2023,luoRecONOnlineLearning2023,li2024f}.
Scene representation instead encodes localized B-mode observations into a
coherent 3D model, with geometry measured, estimated, or jointly refined.
Early approaches mainly used implicit neural representations to learn continuous
spatial-to-intensity or ultrasound-parameter mappings
\cite{yeungSensorlessVolumetricReconstruction2024,
wysockiUltraNeRFNeuralRadiance2024,
dagliNeRFUSRemovingUltrasound2024}.
Although expressive, repeated neural-field evaluations incur substantial
optimization and rendering costs.
More recently, 3D Gaussian Splatting (3DGS) explicitly represents scenes with
anisotropic Gaussian primitives optimized through differentiable rasterization,
enabling efficient rendering and direct inspection
\cite{kerbl3DGaussianSplatting2023}.
These advantages have driven growing interest in Gaussian-based US reconstruction
\cite{eidUltraGaussUltrafastGaussian2025,
yangUltraGSRealTimePhysicallyDecoupled2025,
duelmerUltraGRayPhysicsBasedGaussian2026}.

Nevertheless, sensorless Gaussian-based 3D US reconstruction remains
challenging in both acquisition and pipeline design.
From the acquisition side, sensorless reconstruction lacks measured frame-wise
poses, while limited inter-slice overlap and view-dependent speckle make
unconstrained motion recovery difficult and prone to accumulated geometric
inconsistency.
Existing methods mainly estimate frame-wise or inter-frame 6-DoF probe
transformations from neighboring frames or temporal sequences
\cite{prevostDeepLearningSensorless2017,guoSensorlessFreehand3D2020,
guoUltrasoundVolumeReconstruction2023,luoRecONOnlineLearning2023,li2024f}.
However, this pose-centric formulation underuses the regularity of many
practical sweeps, where neighboring slices are approximately parallel and probe
motion is predominantly translational, introducing redundant degrees of freedom
\cite{huangMMSEReconstruction2008,liTUSREC2024}.
Exploiting this regularity for globally consistent geometry without independent
per-frame poses therefore remains open.
From the representation side, directly transferring the original 3DGS pipeline
to ultrasound is inadequate because both its renderer and Gaussian adaptation
are designed around optical image formation.
Standard 3DGS relies on perspective projection and front-to-back alpha
compositing, whereas B-mode ultrasound arises from acoustic propagation,
attenuation, reflection, and finite-slice scattering, making image formation
depend on both localized echo generation and tissues along the preceding
propagation path.
This mismatch demands an ultrasound-specific renderer that jointly models local
echo formation and path-dependent transmission.
Moreover, Gaussian adaptation must also be reconsidered: image-space gradients
can be corrupted by speckle, while opacity cannot capture path-dependent
primitive utility.
Therefore, robustly estimating per-Gaussian contributions for adaptive
densification and pruning remains another key challenge.

To address these challenges, we propose \methodname{}, a Gaussian-based
framework for sensorless 3D US reconstruction that integrates compact geometry
modeling, ultrasound-specific rendering, and a tailored adaptation strategy,
enabling high-fidelity volumetric reconstruction, efficient optimization, and
high-speed rendering without external tracking.
To establish globally consistent sensorless geometry, \methodname{} employs an
\emph{Acquisition-constrained Frame Geometry Model (AFGM)}, which replaces
independent frame-wise pose recovery with a compact acquisition model, providing
a stable and coherent spatial scaffold for Gaussian reconstruction under
regular, translation-dominant sweeps.
Given this geometry, we develop a \emph{Physics-aware Local-Global Renderer
(PLGR)} as the core scene-representation component of \methodname{}, where a
finite-slice local pathway models localized echo formation and a complementary
global pathway captures path-dependent transmission from causally preceding
tissues, preserving slice-specific structural detail and propagation-consistent
acoustic context to improve B-mode reconstruction fidelity.
Building directly on PLGR, a \emph{Contribution-guided Gaussian Adaptation
Strategy (CGAS)} uses residual-weighted local contributions to densify
under-reconstructed regions and joint local-global contributions for pruning,
concentrating Gaussian capacity on insufficiently modeled structures while
preserving primitives with meaningful path effects and aligning population
control with ultrasound image formation.
Comprehensive experiments on four synthetic and real-world datasets covering
liver, spine, arm, and breast anatomy demonstrate robust reconstruction across
diverse acquisition settings, underscoring the potential of \methodname{} for
efficient and practical sensorless 3D US reconstruction in real-world
volumetric ultrasound applications.
